\documentclass[11pt]{article}

\usepackage[margin=1in]{geometry}
\usepackage{amsmath, amssymb, amsfonts}
\usepackage{enumitem}
\usepackage{hyperref}

\title{Corpus Generation and Search}
\author{kinematic-classifier-sandbox}
\date{\today}

\begin{document}

\maketitle

\section{Purpose}

This document explains the corpus-side methodology modules: how witness
trajectories are parameterized, how candidate corpora are searched and scored,
and how study candidates are promoted from those corpus surfaces.

\section{Problem And Notation}

The corpus-side methodology is trying to solve:
\begin{equation}
    \text{How do we generate and select datasets that are informative enough to test classifiers without making the task accidentally trivial or biased?}
\end{equation}

The most important objects are:
\begin{itemize}[leftmargin=1.5em]
    \item \(\tau\): one generated trajectory,
    \item \(\mathcal{D}\): a corpus candidate,
    \item \(\theta_{\text{class}}\): class-specific motion parameters,
    \item \(\theta_{\text{tier}}\): difficulty-tier controls,
    \item \(\theta_{\text{noise}}\): corruption parameters,
    \item \(\theta_{\text{sampling}}\): timing and sample-count parameters,
    \item \(S_k\): scalar score for corpus candidate \(k\),
    \item \(\mathbf{o}_k\): Pareto objective vector for candidate \(k\).
\end{itemize}

\section{Trajectory Parameterization And Witness Problems}

The base synthetic witness-problem logic lives in:
\begin{itemize}[leftmargin=1.5em]
    \item \texttt{trajectory\_generator.py}
    \item \texttt{corpus\_objectives.py}
\end{itemize}

The trajectory generator is the repo's explicit map from class semantics to
parameterized synthetic observations. In abstract form, each generated
trajectory is
\begin{equation}
    \tau = \tau(\theta_{\text{class}}, \theta_{\text{tier}}, \theta_{\text{noise}}, \theta_{\text{sampling}}, \xi),
\end{equation}
where the class parameters encode the nominal motion family, the tier
parameters encode difficulty and corruption, and \(\xi\) is random seed
realization.

This is why the corpus layer matters mathematically: later classifier and
adequacy studies are only as meaningful as the induced distribution over
\(\tau\).

\texttt{corpus\_objectives.py} provides the declarative objective layer that
tells the generator what kind of corpus pressure to create, including class-pair
targets, feature-excitation targets, leakage constraints, and runtime budgets.

\section{Adaptive Stress And Environment-Aware Corpus Layers}

The current corpus-shaping surfaces include:
\begin{itemize}[leftmargin=1.5em]
    \item \texttt{adaptive\_stress\_corpus.py}
    \item \texttt{environment\_aware\_corpus.py}
    \item \texttt{quality\_diversity\_corpus.py}
    \item \texttt{objective\_driven\_qd\_archive.py}
\end{itemize}

These modules vary difficulty, environment assumptions, and stress conditions
rather than treating one synthetic corpus as fixed forever. Their shared role is
to perturb the corpus distribution along interpretable axes:
\begin{equation}
    \mathcal{D}
    \rightarrow
    \mathcal{D}'(\text{noise}, \text{irregularity}, \text{outliers}, \text{stress}).
\end{equation}

\section{Corpus Autodevelopment And Pareto Selection}

The automated corpus-development layer is implemented through:
\begin{itemize}[leftmargin=1.5em]
    \item \texttt{corpus\_autodevelopment.py}
    \item \texttt{corpus\_gym.py}
\end{itemize}

This layer treats corpus design as an optimization problem over:
\begin{itemize}[leftmargin=1.5em]
    \item balance,
    \item feature excitation,
    \item class-pair boundary coverage,
    \item leakage penalties,
    \item adequacy score.
\end{itemize}

Candidate \(k\) receives a score of the form
\begin{equation}
    S_k
    =
    w_b B_k
    + w_c C_k
    + w_f F_k
    + w_d D_k
    - w_\ell L_k
    - w_t T_k
    - w_g G_k,
\end{equation}
where:
\begin{itemize}[leftmargin=1.5em]
    \item \(B_k\) is balance score,
    \item \(C_k\) is boundary coverage score,
    \item \(F_k\) is feature excitation score,
    \item \(D_k\) is difficulty diversity score,
    \item \(L_k\) is leakage penalty,
    \item \(T_k\) is triviality penalty,
    \item \(G_k\) is degeneracy penalty.
\end{itemize}
That is the logic implemented across
\texttt{\_feature\_excitation\_score},
\texttt{\_boundary\_coverage\_score},
\texttt{\_balance\_score},
\texttt{\_leakage\_penalty},
\texttt{\_triviality\_penalty}, and
\texttt{\_degeneracy\_penalty}.
The scalar score is useful for selection, but the accompanying Pareto surface
ensures that non-dominated tradeoffs are not erased.

The same module also computes a Pareto objective vector
\begin{equation}
    \mathbf{o}_k =
    \big[
        B_k,\,
        C_k,\,
        F_k,\,
        D_k,\,
        -L_k,\,
        -T_k,\,
        -G_k
    \big],
\end{equation}
and marks a candidate as dominated when another candidate is no worse in every
coordinate and strictly better in at least one.

The artifact
\texttt{artifacts/corpus\_autodevelopment\_v1/corpus\_autodevelopment\_numeric\_walkthrough.md}
now works through one real selected candidate numerically. It substitutes the
implemented score terms into the scalar objective, expands the
difficulty-distribution subscore against the configured target fractions, and
then compares the selected candidate against the highest-scoring rejected one.
That artifact is the concrete bridge from the symbolic objective above to
\texttt{render\_corpus\_autodevelopment\_numeric\_walkthrough\_markdown(...)}.

\section{Assumptions Behind Corpus Selection}

The corpus autodevelopment layer assumes:
\begin{itemize}[leftmargin=1.5em]
    \item positive adequacy terms should be maximized,
    \item leakage, triviality, and degeneracy terms should be penalized,
    \item a single scalar score is useful for default selection,
    \item but Pareto non-dominance should still be preserved for auditability.
\end{itemize}

\section{Corpus Search And Baseline Ranking}

The broader search/comparison layer is implemented through:
\begin{itemize}[leftmargin=1.5em]
    \item \texttt{corpus\_search\_baseline.py}
    \item \texttt{corpus\_synthesis\_comparison.py}
    \item \texttt{generic\_corpus\_exploration.py}
    \item \texttt{selected\_generated\_corpus.py}
\end{itemize}

These modules are the ``search shell'' around the witness-problem generator and
adequacy machinery. The core methodological claim is:
\begin{equation}
    \text{corpus search}
    \neq
    \text{randomly generate more trajectories}.
\end{equation}
Instead, the repo is moving toward objective-driven search over corpus
properties that matter to identifiability and evaluation credibility.

\subsection*{Generic Corpus Explorer}

\texttt{generic\_corpus\_exploration.py} is the clearest implementation of the
repo's Corpus Explorer idea. It starts from a heterogeneous pool of
backend-specific candidate trajectories and scores each executed run by a mixed
utility rather than by one raw classifier metric:
\begin{equation}
    U_{\text{explore}}
    =
    0.22 \cdot \text{validity}
    + 0.18 \cdot \text{coverage novelty}
    + 0.18 \cdot \text{boundary score}
    + 0.18 \cdot \text{classifier stress}
    + 0.12 \cdot \text{environment score}
    + 0.12 \cdot \text{provenance completeness}.
\end{equation}
This is a deliberately mixed objective: it rewards not only hard or informative
examples, but also archive-cell novelty, environment structure, and audit
traceability.

\subsection*{Explorer Archive Logic}

The explorer also constructs archive cells over backend, scenario family,
target class, and difficulty tier, then preserves the elite with maximal
\texttt{total\_utility} in each cell. The selected corpus is judged against a
same-size random baseline by coverage:
\begin{equation}
    \Delta_{\text{coverage}}
    =
    \#\{\text{selected archive cells}\}
    -
    \#\{\text{random-baseline archive cells}\}.
\end{equation}
That is why the explorer can claim more than ``high utility'': it can ask
whether its selected set covers more meaningful behavioral cells than a naïve
random draw.

\subsection*{Worked Example}

The artifact
\texttt{artifacts/generic\_corpus\_exploration/generic\_corpus\_explorer\_numeric\_walkthrough.md}
now expands one selected corpus row into its utility decomposition, archive-cell
role, elite interpretation, and selected-versus-random coverage comparison.
That artifact is the Explorer-side proof that corpus selection is not only a
visual dashboard but a numeric utility and coverage argument on one concrete
selected row.

\section{Corpus Gym}

\texttt{corpus\_gym.py} reframes corpus generation as a targeted environment.
Its question is:
\begin{equation}
    \text{Can we specify desired failure pressure or feature geometry and reward trajectories that approach it?}
\end{equation}

The main objects are:
\begin{itemize}[leftmargin=1.5em]
    \item a target specification,
    \item an action that perturbs tier, measurement, irregularity, outlier, or
    step settings,
    \item a structured reward decomposition,
    \item and an episode object containing target, action, trajectory,
    diagnostics, and reward.
\end{itemize}

The reward in \texttt{CorpusGymReward} is multi-term:
\begin{equation}
    U_{\text{gym}}
    =
    0.22 \cdot V
    + 0.14 \cdot E
    + 0.14 \cdot G
    + 0.14 \cdot B
    + 0.14 \cdot S
    + 0.12 \cdot P
    - 0.10 \cdot L
    - 0.14 \cdot I,
\end{equation}
where \(V\) is class validity, \(E\) is feature excitation match, \(G\) is
coverage gain, \(B\) is boundary closeness, \(S\) is classifier-stress
pressure, \(P\) is prior-sensitivity pressure, \(L\) is leakage penalty, and
\(I\) is physical invalidity penalty. The component terms are computed through
\texttt{\_class\_validity\_score(...)},
\texttt{\_feature\_excitation\_score(...)},
\texttt{\_coverage\_gain\_score(...)},
\texttt{\_boundary\_closeness\_score(...)},
\texttt{\_classifier\_stress\_score(...)},
\texttt{\_prior\_sensitivity\_score(...)},
\texttt{\_leakage\_penalty(...)}, and
\texttt{\_physical\_invalidity\_penalty(...)}.

\subsection*{Coverage, Leakage, and Invalidity Terms}

Three Gym subterms are especially important because they stop the search from
rewarding obviously biased or pathological cases. The implemented coverage-gain
term is
\begin{equation}
    G
    =
    0.40 \cdot \text{tier match}
    + 0.30 \cdot \text{class match}
    + 0.30 \cdot \text{novelty}(a),
\end{equation}
where the novelty term averages capped measurement, irregularity, outlier, and
step scales induced by the action. The leakage penalty is
\begin{equation}
    L
    =
    0.30 \cdot \text{duration risk}
    + 0.30 \cdot \text{sample-count risk}
    + 0.40 \cdot \text{noise risk},
\end{equation}
which prevents the Gym from rewarding trajectories merely because duration,
sample count, or noise level become class-identifying shortcuts. The physical
invalidity penalty guards against non-increasing time grids and implausibly
large absolute accelerations.

\subsection*{Environment Contract}

\texttt{CorpusGymEnvironment} makes the search loop explicit. It provides
\texttt{reset(target)}, \texttt{simulate(action)}, \texttt{step(action)},
\texttt{trajectory()}, \texttt{score(trajectory)}, and
\texttt{render\_diagnostics(trajectory)}. Methodologically, this is
\begin{equation}
    \text{target}
    \xrightarrow{\text{reset}}
    \text{state}
    \xrightarrow{\text{action}}
    \text{trajectory, reward, diagnostics}.
\end{equation}
That is why the Gym should be read as a reusable corpus-search interface rather
than a one-off artifact helper.

\subsection*{Objective-To-Gym Execution Bridge}

\texttt{objective\_corpus\_gym\_runner.py} connects the declarative
corpus-objective layer to the Gym layer. It defines
\begin{equation}
    \Psi_{\text{target}}(\text{objective}) \rightarrow \text{CorpusGymTarget},
    \qquad
    \Psi_{\text{action}}(\text{candidate}) \rightarrow \text{CorpusGymAction},
\end{equation}
and then executes
\begin{equation}
    \text{objective}
    \rightarrow
    \text{candidate}
    \rightarrow
    \text{target, action}
    \rightarrow
    \text{episode}
    \rightarrow
    \text{validated trajectory run}.
\end{equation}
This is the critical bridge from declarative study design to executed explorer
or Gym records. Without it, the Gym would remain an isolated search surface.

\subsection*{Worked Example}

The artifact
\texttt{artifacts/corpus\_gym/corpus\_gym\_numeric\_walkthrough.md}
now works through one real Gym episode numerically. It records the selected
target, the action scales, the resulting trajectory diagnostics, and the
implemented reward substitution term by term. That artifact is the Gym-side
proof that the search reward is not only qualitative prose but a concrete score
decomposition on one executed trajectory.

\section{Objective-Driven Quality-Diversity Archive}

\texttt{objective\_driven\_qd\_archive.py} preserves diverse elites instead of
only selecting one scalar winner. For one successful archive cell, the current
elite utility is
\begin{equation}
    U_{\text{archive}}
    =
    0.30 \cdot \text{validity}
    + 0.25 \cdot \text{acceleration-range pressure}
    + 0.25 \cdot \text{classifier stress}
    + 0.20 \cdot (1 - \text{mean margin}).
\end{equation}
The cell identity itself depends on discretized buckets over class, tier,
backend, duration, acceleration range, entropy, and prior-flip threshold. This
means the archive is tracking both elite quality and structural coverage.

\section{Study Candidate Generation And Search Surfaces}

The study proposal layer is implemented through:
\begin{itemize}[leftmargin=1.5em]
    \item \texttt{candidate\_generation.py}
    \item \texttt{study\_candidate\_generation.py}
    \item \texttt{study\_candidate\_protocol.py}
    \item \texttt{capability\_aware\_search.py}
\end{itemize}

These modules connect corpus generation to feature/classifier proposal logic.
\texttt{candidate\_generation.py} is the lower-level sampler layer underneath
that study-selection logic. It mutates and combines backend-compatible candidate
specifications before the higher-level protocol decides which ones should be
promoted.
For a candidate study built from class pair \(p\), classifier \(m\), feature set
\(f\), and prior regime \(r\), the current static score is effectively
\begin{align}
    Q_{\text{static}}(p,m,f,r)
    &=
    a_1 \cdot \text{feature-class compatibility}
    + a_2 \cdot \text{expected separability}
    + a_3 \cdot \text{classifier fit} \notag \\
    &\quad
    + a_4 \cdot \text{corpus coverage}
    + a_5 \cdot \text{dimensional transfer}
    + a_6 \cdot \text{implementation readiness} \notag \\
    &\quad
    + a_7 \cdot (1 - \text{dependency risk})
    - a_8 \cdot \text{cumulative-history risk}
    - a_9 \cdot \text{prior-sensitivity risk}.
\end{align}
The current implementation in \texttt{study\_candidate\_generation.py} uses the
explicit weights
\[
(a_1,a_2,a_3,a_4,a_5,a_6,a_7,a_8,a_9)
=
(0.18,0.18,0.14,0.14,0.12,0.12,0.12,0.10,0.10).
\]
This is why study candidate generation is a methodology module rather than just
a report writer: it is already combining feature, corpus, classifier, and prior
information into a promotion surface.

In the current code, that high-level score is instantiated through subterms such
as \texttt{feature\_class\_compatibility\_score},
\texttt{expected\_separability\_score},
\texttt{classifier\_assumption\_fit},
\texttt{corpus\_coverage\_score},
\texttt{dimensional\_transfer\_score},
\texttt{implementation\_readiness\_score},
\texttt{feature\_dependency\_risk},
\texttt{cumulative\_double\_counting\_risk}, and
\texttt{prior\_sensitivity\_risk}.
These are also the actual score columns written to
\texttt{static\_candidate\_scores.csv}.

\subsection*{Assumptions}

The static score assumes that feature/class compatibility and oracle
separability should dominate early screening, that corpus coverage and
classifier-family fit are second-order but still important, and that 3D
transferability plus implementation readiness should matter before the repo is
fully generalized. It also assumes a two-stage selection process:
\begin{equation}
    \text{static screening}
    \rightarrow
    \text{Monte Carlo confirmation}
    \rightarrow
    \text{promotion decision}.
\end{equation}

\subsection*{Monte Carlo Confirmation Layer}

After static screening, the same module uses common-harness benchmark outputs to
compute
\begin{equation}
    Q_{\text{mc}}
    =
    0.60 \cdot \text{accuracy}
    + 0.25 \cdot (1 - \text{prior flip fraction})
    + 0.15 \cdot (1 - \max(0, \text{oracle gap})).
\end{equation}
The current promote gate is approximately
\begin{align}
    \text{compatible}
    &\land
    Q_{\text{static}} \ge 0.45
    \land
    Q_{\text{mc}} \ge 0.90
    \notag \\
    &\land
    \text{accuracy} \ge 0.83
    \land
    \text{prior flip fraction} \le 0.12.
\end{align}
If these conditions are not met but the candidate still has acceptable
compatibility and classifier performance, the decision typically drops to
\texttt{revise}; otherwise it falls to \texttt{reject} or \texttt{defer}. That
is the real evidence gate between proposal generation and study promotion.

\section{Candidate Generation As A Sampler Family}

\texttt{candidate\_generation.py} does not emit one candidate per objective. For
objective \(o\) and backend \(b\), it effectively builds
\begin{equation}
    \mathcal{C}(o,b)
    =
    \mathcal{C}_{\text{random}}
    \cup
    \mathcal{C}_{\text{grid}}
    \cup
    \mathcal{C}_{\text{lhs}}
    \cup
    \mathcal{C}_{\text{boundary mutation}}
    \cup
    \mathcal{C}_{\text{archive mutation}}
    \cup
    \mathcal{C}_{\text{stress mutation}}.
\end{equation}
This is why the module is already a search policy rather than a single
constructor.
In probabilistic language it behaves like a mixture proposal family
\begin{equation}
    q(c \mid o,b)
    =
    \sum_s \pi_s q_s(c \mid o,b),
\end{equation}
where \(s\) ranges over sampler families and the mixture weights are realized
through fixed per-family generation budgets rather than learned online.

\section{Capability-Aware Search Planning}

\texttt{capability\_aware\_search.py} maps backend capabilities into search
policies. At the methodology level, the planner can be viewed as
\begin{equation}
    \Pi(\text{backend capabilities})
    \rightarrow
    \{\text{recommended search methods}, \text{budget class}, \text{planner rationale}\}.
\end{equation}
The planner conditions on runtime class, dimensionality, environment support,
sequential-control support, and stochastic versus deterministic execution.
More concretely, if \(\kappa\) denotes the backend capability vector, the rule
set behaves like
\begin{equation}
    \Pi(\kappa)
    =
    \big(
        M_{\text{runtime}}(\kappa),
        M_{\text{environment}}(\kappa),
        M_{\text{control}}(\kappa),
        M_{\text{stochastic}}(\kappa)
    \big),
\end{equation}
followed by deduplication into one backend plan row. Cheap stochastic backends
receive broad search budgets, while expensive deterministic backends are pushed
toward smaller DOE, surrogate-assisted search, and cache-priority execution.

\section{Study Candidate Protocol And Validation Ladder}

\texttt{study\_candidate\_protocol.py} defines the contract under the
generated-candidate story. It specifies the \texttt{StudyCandidate} schema, the
\texttt{ValidationLadder} schema, and the terminal decision vocabulary
\{\texttt{promote}, \texttt{revise}, \texttt{reject}, \texttt{defer}\}. The
promotion story is therefore
\begin{equation}
    \text{candidate specification}
    \rightarrow
    \text{validation ladder evidence}
    \rightarrow
    \text{terminal decision in a constrained vocabulary}.
\end{equation}
The protocol also defines ordered validation levels, including static
compatibility, corpus adequacy, feature separability, oracle separability,
classifier performance, posterior and calibration quality, prior sensitivity,
stress robustness, dimensional transfer assessment, and final promotion
decision. In compact form:
\begin{equation}
    \text{proposal}
    \rightarrow
    \{\ell_1,\ell_2,\dots,\ell_{10}\}
    \rightarrow
    d,
    \qquad
    d \in \{\text{promote}, \text{revise}, \text{reject}, \text{defer}\}.
\end{equation}
This is what makes the study layer auditable rather than ad hoc.

\section{Generated Class And Feature Exploration}

\texttt{generated\_corpus\_features.py} routes objective-driven generated
trajectories back through the real feature pipeline. The methodological flow is
\begin{equation}
    \text{generated candidate}
    \rightarrow
    \text{executed trajectory}
    \rightarrow
    \text{validity-adjusted label}
    \rightarrow
    \text{feature row}
    \rightarrow
    \text{excitation and separability analysis}.
\end{equation}

\texttt{corpus\_classifier\_scoring.py} then evaluates how the current
classifier ladder behaves on those generated trajectories. Its stress proxy is
approximately
\begin{equation}
    \text{stress}
    \approx
    0.5 \cdot (1 - \text{margin})
    + 0.5 \cdot \text{entropy}
    + 0.35 \cdot \mathbf{1}\{\text{final prediction wrong}\}.
\end{equation}
This is what turns generated-corpus exploration into an interpretable source of
classifier pressure rather than just another batch of synthetic rows.

\section{Methodological Flow}

The important transition is:
\begin{align}
    \text{trajectory generator}
    &\rightarrow
    \text{corpus candidate}
    \rightarrow
    \text{adequacy-scored corpus}
    \notag \\
    &\rightarrow
    \text{backend-aware candidate population}
    \rightarrow
    \text{study candidate}
    \rightarrow
    \text{validation ladder}
    \rightarrow
    \text{promotion / revise / reject / defer}.
\end{align}

The critical point is that synthetic witness problems are not the end product.
They are reusable inputs to a study-selection and methodology-audit loop.

\section{Artifact Map}

Primary artifacts for this layer:
\begin{itemize}[leftmargin=1.5em]
    \item \texttt{artifacts/corpus\_autodevelopment\_v1/}
    \item \texttt{artifacts/corpus\_synthesis\_comparison/}
    \item \texttt{artifacts/study\_candidate\_generation/}
    \item \texttt{artifacts/adaptive\_stress\_corpus/}
    \item \texttt{artifacts/environment\_aware\_corpus/}
\end{itemize}

\end{document}
